\begin{frame}{Апробация метода начнётся после модельных проверок}
\SlideTableTitle{Публикации по теме исследования}
{\fontsize{12}{14.4}\selectfont
\begin{tabularx}{\linewidth}{|Y|L{0.30\linewidth}|L{0.17\linewidth}|}
\hline Работа & Площадка & Состояние \\\hline
Методы оценки болевого состояния по биомедицинским данным: целевые переменные, доступные наборы данных и требования к валидации & Конференция «Микроэлектронные имплантируемые нейроинтерфейсы 2026» & Направлена \\\hline
Вычислительные модели нейронных систем в нейроинженерии: источники ограничений, данные и валидация & Журнал «Биомедицинская радиоэлектроника» & Рукопись, ноябрь 2026 \\\hline
Коннектомно-ограниченное предобучение для биофизического моделирования ECAP при SCS: предпосылки и критерии проверки & Журнал «Биомедицинская радиоэлектроника» & Рукопись, март 2027 \\\hline
\end{tabularx}\par}
\SlideGap
Индивидуальный план предусматривает 16 работ в 2026--2030 годах: 11 журнальных и 5 конференционных. Обзоры обосновывают постановку задачи и не заменяют апробацию модельных результатов.
\end{frame}
